% TOP_PROBLEM: {"id": 2301017, "title": "Research Problems in Function Theory — Problem 1.17", "category_id": 8, "category": {"id": 8, "name": "analysis", "display_name": "Analysis", "description": "Limits, continuity, calculus, and function theory.", "slug": "analysis", "order_index": 8, "created_at": "2026-07-31T15:26:25.671Z"}, "status": "partially_solved"}
% FIRST_POSTED: null
% ENTRY_KIND: statement_only
% Imported from: batch10.tex; UnsolvedMath ID 2301017
\documentclass[11pt]{article}
\usepackage[margin=25mm]{geometry}
\usepackage{fontspec}
\setmainfont{Latin Modern Roman}
\usepackage{amsmath,amssymb,mathtools}
\usepackage{unicode-math}
\setmathfont{Latin Modern Math}
\usepackage{xurl,hyperref}
\hypersetup{colorlinks=true,urlcolor=blue}
\setcounter{secnumdepth}{0}
\setlength{\parindent}{0pt}
\setlength{\parskip}{5pt}
\setlength{\emergencystretch}{3em}
\newcommand{\sref}[2]{\hyperlink{source:#1:#2}{[S#2]}}
\newcommand{\cataloguescope}[1]{\paragraph{Recorded target.}#1}
\begin{document}
% UnsolvedMath ID 2301017; source catalogue code AMR-022-1017
\setcounter{section}{2301016}
\section{Research Problems in Function Theory — Problem 1.17}\label{2301017}
{\small\sffamily UnsolvedMath ID 2301017 · Analysis}
\subsection{Definitions and mathematical statement}

\paragraph{Shared notation.}$\mathbb N=\{1,2,\ldots\}$, and $\mathbb Z,\mathbb Q,\mathbb R,\mathbb C$ denote the integers, rationals, reals and complex numbers. Any other conventions are specified locally.


For a nonconstant meromorphic function $f$ on $\mathbb C$, let $n(r,a,f)$ count the solutions of $f(z)=a$ in $|z|\le r$, with multiplicity; $a=\infty$ counts poles. Let $\bar n$ count each point once. The integrated count is
\[
 N(r,a,f)=n(0,a,f)\log r+\sum_{0<|z|\le r,\ f(z)=a}\log(r/|z|),
\]
where the sum repeats each point according to multiplicity; $\bar N$ has the same formula without repetitions. Write $\log^+x=\max(\log x,0)$ and
\[
 m(r,\infty,f)=\frac1{2\pi}\int_0^{2\pi}\log^+|f(re^{it})|\,dt,\quad
 m(r,a,f)=\frac1{2\pi}\int_0^{2\pi}\log^+\frac1{|f(re^{it})-a|}\,dt\quad(a\in\mathbb C).
\]
The Nevanlinna characteristic is $T(r,f)=m(r,\infty,f)+N(r,\infty,f)$. Its order and lower order are
\[
 \rho(f)=\limsup_{r\to\infty}\frac{\log T(r,f)}{\log r},\qquad
 \lambda(f)=\liminf_{r\to\infty}\frac{\log T(r,f)}{\log r}.
\]
The deficiency is $\delta(a,f)=1-\limsup_{r\to\infty}N(r,a,f)/T(r,f)$, and the reduced-count index is $\Theta(a,f)=1-\limsup_{r\to\infty}\bar N(r,a,f)/T(r,f)$. A value is deficient if $\delta(a,f)>0$. An entire function is holomorphic on all of $\mathbb C$; for an entire function, $M(r,f)=\max_{|z|=r}|f(z)|$.
\paragraph{Statement recorded in the supplied source.}
For $\rho>0$ let $\mathcal E_\rho$ be the entire functions of order exactly $\rho$, and define
\[
 C(\rho)=\sup_{f\in\mathcal E_\rho}\liminf_{r\to\infty}\frac{\log M(r,f)}{T(r,f)}.
\]
Is $C(\rho)=\pi\rho$ for every $\rho>1/2$? For comparison, the source states $C(\rho)=\pi\rho/\sin(\pi\rho)$ when $0<\rho<1/2$. \sref{2301017}{2}
\subsection{Short English statement}
Determine the sharp upper bound for this lower limiting maximum-modulus ratio, specifically the proposed value $\pi$ times the order above one-half. The source records Govorov’s proof for entire functions. \sref{2301017}{2}
\subsection{Sources}
\begin{description}
\item[\hypertarget{source:2301017:1}{S1}] ULAM.AI and the UnsolvedMath Contributors, UnsolvedMath: A Curated Collection of Open Mathematics Problems, record 2301017 (AMR-022-1017). CC BY 4.0. \url{https://huggingface.co/datasets/ulamai/UnsolvedMath}.
\item[\hypertarget{source:2301017:2}{S2}] Walter K. Hayman and Eleanor F. Lingham, Research Problems in Function Theory, Fiftieth Anniversary Edition (2018), arXiv:1809.07200. Problem 1.17, its full Update 1.17, and chapter notation; original author TeX. \url{https://arxiv.org/abs/1809.07200}.
\end{description}
\label{end:2301017}
\end{document}
